
\subsubsection{6.1 Format-Dependency}

The 0.24 AUROC gap between max probability and semantic entropy results from format mismatch. Semantic entropy assumes responses contain meaningful semantic content for NLI comparison. Single-letter multiple-choice answers violate this assumption.

Semantic entropy was developed and validated on free-form question answering where responses are sentences or paragraphs. The method's design is reasonable for that setting. The present results indicate it does not transfer to multiple-choice format without modification.

\subsubsection{6.2 Practical Implications}

For multiple-choice format hallucination detection, simple confidence-based methods are sufficient. The additional computational cost of multi-sample generation and NLI-based clustering does not improve discrimination on this format.

Method selection should consider task format. More sophisticated methods do not uniformly outperform simpler alternatives.

\subsubsection{6.3 Limitations}

\textbf{Sample size.} Experiments used 50 questions (proof-of-concept level). Full validation requires the complete 817-sample TruthfulQA dataset. The relative rankings (max probability > choice entropy > semantic entropy) are expected to hold, but exact AUROC values have associated uncertainty.

\textbf{Single model.} Only Llama-3-8B-Instruct was evaluated. Results may not generalize across architectures or model sizes.

\textbf{Multiple-choice format only.} Free-form generation was not evaluated. Semantic entropy may perform differently on tasks where responses contain substantial semantic content.

\textbf{Sample count.} Semantic entropy used 5 samples per question. The original work used 5--10 samples. Additional samples might improve performance, though this would not address the fundamental format mismatch.

\subsubsection{6.4 Connection to Prior Work}

The finding that semantic entropy underperforms on multiple-choice format does not contradict Kuhn et al. [2023], who evaluated on free-form question answering. The present results extend understanding of the method's scope conditions: semantic entropy requires responses with semantic content suitable for NLI comparison.

The finding that max probability achieves strong discrimination (AUROC = 0.81) aligns with Kadavath et al. [2022], who demonstrated that LLM confidence signals correlate with correctness.

